\scanpage{27}

Donc divisant $\underline{\mathcal O}_{X,x}$ par un idéal défini par $\deleted{\unsure{k}}n$ éléments, on \deleted{\unsure{obtient|de ce fait}} trouve par le th général \deleted{\unsure{un}} (\unsure{appliquer|appliqué} : $\underline{\mathcal O}_{\unsure{X|x},x}$ intègre formellt équidim. de dim $\geq K+1=(n+k)+1$] un anneau qui est connexe en dim $\geq\deleted{(n+k)}\inserted{K-n}=k$, cqfd.

\uline{Cor. 1} au th. 2 Variante avec \deleted{composantes} ``\unsure{géométriquemt|géométriqt} \deleted{\unsure{int.|iréd.}} unibranche'', \unsure{*|``} \deleted{\unsure{et}} fibres géométriques.

\uline{Cor. 2} Variante avec $f$ propre univ. ouvert.

\uline{Cor. 3} $X\subset P^r_k$, \deleted{$X$ \unsure{de di}} \inserted{les composantes} de $\unsure{\alpha|X}$ de dim $\geq k+1$, et $X$ connexe en dim $\geq k$. Alors pour tt système de $\ell$ hyperplans $H_1\text{--}H_\ell$, \inserted{\unsure{le schéma|la section}} \deleted{\unsure{$Z$}} $X\cap H_1\cap-\cap H_\ell=Y$ \unsure{satisfait} les conditions analogues avec \unsure{au|$\ell$ au} lieu de $k$.
% slanted
\margin{Généraliser au cas d'un morphisme $X\longrightarrow P_k^r$.}
% slanted
\margin{Peut-on généraliser de façon analogue le th. local ? Facile, en utilisant le thé. \unsure{tel} \unsure{en raisonnant|correspondant} \unsure{relatif}.}

% archive sheet number: 26
